\section*{Capítulo \ref{ch:b1:second-law-entropy} --- El segundo principio: la entropía}

\begin{solution}{exo:b1:second-law-entropy:1}
$\Delta S = mL/T$: fusión $334000/273 = \qty{1.22}{kJ/K}$; ebullición $2260000/
373 = \qty{6.06}{kJ/K}$. La vaporización libera del todo a las moléculas
unas de otras --- mucho más desorden microscópico y mucho más calor por
\omterm{def:b1:kinetic-theory:temperature}{kelvin}.
\end{solution}

\begin{solution}{exo:b1:second-law-entropy:2}
$\Delta S = mc\ln(T_2/T_1) = 4180\ln(353/293) = \qty{0.78}{kJ/K}$ --- una función
de estado: independiente del método (la entropía \emph{creada} no lo es).
\end{solution}

\begin{solution}{exo:b1:second-law-entropy:3}
$\Delta S = nR\ln2 = \qty{5.8}{J/K}$. Isoterma reversible: $Q = nRT\ln2$
del baño a $T$, $S_{\mathrm{exch}} = nR\ln2$, creada $0$. Joule:
$Q = 0$, $S_{\mathrm{exch}} = 0$, creada \qty{5.8}{J/K}.
\end{solution}

\begin{solution}{exo:b1:second-law-entropy:4}
$S = nC_{V,m}\ln T + nR\ln V$ constante: $T^{C_{V,m}}V^R$ constante, es decir
$TV^{R/C_{V,m}} = TV^{\gamma - 1}$ constante. Compresión brusca: $Q = 0$, luego
$S_{\mathrm{exch}} = 0$ y $\Delta S = S_{\mathrm{created}}$; con $T$: $300 \to
643$ K y $V$: $24.9 \to 10.7$ L: $\Delta S = 20.8\ln(643/300) + 8.314\ln(10.7/
24.9) = 15.9 - 7.0 = \qty{+8.9}{J/K}$ creados.
\end{solution}

\begin{solution}{exo:b1:second-law-entropy:5}
Bloque: $450\ln(293/373) = \qty{-109}{J/K}$; lago: $+36000/293 = \qty{+123}{J/K}$;
universo: $+\qty{14}{J/K}$.
\end{solution}

\begin{solution}{exo:b1:second-law-entropy:6}
$T_f = \qty{350}{K}$; $S = 4180[\ln(350/400) + \ln(350/300)] = 4180 \times
0.0206 = \qty{86}{J/K}$, todo creado. $(T_1 + T_2)^2 - 4T_1T_2 = (T_1 - T_2)^2
\geq 0$, luego el argumento del logaritmo es $\geq 1$.
\end{solution}

\begin{solution}{exo:b1:second-law-entropy:7}
$S = Q(1/T_c - 1/T_h)$: $1000(1/280 - 1/300) = \qty{0.24}{J/K}$; del horno a
la habitación: $1000(1/300 - 1/1000) = \qty{2.3}{J/K}$, diez veces más: cuanto
mayor es el salto de temperatura que atraviesa un flujo de calor, más
irreversible --- en la caída del horno a la habitación se derrocha la energía noble.
\end{solution}

\begin{solution}{exo:b1:second-law-entropy:8}
Ciclo: $\Delta S = 0 = Q/T_0 + S_{\mathrm{created}}$ con $S_{\mathrm{created}} \geq 0$,
luego $Q \leq 0$ y $W = -Q \geq 0$: la máquina no puede entregar trabajo. El
mar es un solo termostato: su calor no puede convertirse en trabajo sin un
cuerpo más frío al que verter la entropía.
\end{solution}

\begin{solution}{exo:b1:second-law-entropy:9}
Cada gas se expande de $V$ a $2V$ a $T$ constante: $\Delta S = 2 \times R\ln2
= \qty{11.5}{J/K}$, creados (no hay intercambio de calor, la caja está
aislada). Dos mitades de nitrógeno: retirar el tabique no cambia nada
macroscópico --- $\Delta S = 0$ (paradoja de Gibbs: moléculas idénticas no se ``mezclan'').
\end{solution}

\begin{solution}{exo:b1:second-law-entropy:10}
$Q = \qty{3.6e5}{J}$ a la habitación a \qty{293}{K}: $S = \qty{1230}{J/K}$
creados --- en el filamento (trabajo eléctrico convertido en calor) y en el
paso de ese calor del filamento caliente a la habitación fría; el propio
trabajo eléctrico no traía entropía.
\end{solution}

\begin{solution}{exo:b1:second-law-entropy:11}
$\Delta S_{\text{cuerpo}} = C\ln(T_0/T_1)$ en todos los casos. (a) intercambiada
$C(T_0 - T_1)/T_0$: creada $C[T_1/T_0 - 1 - \ln(T_1/T_0)]$. (b) dos etapas
con razón $r = \sqrt{T_1/T_0}$ cada una: creada $2C[r - 1 - \ln r]$, menor
(para $T_1/T_0 = 2$: $0.193C$ frente a $0.307C$). (c) infinitas etapas:
cada una crea $C[(1 + \epsilon) - 1 - \ln(1 + \epsilon)] \approx C\epsilon^2/2$,
cuya suma tiende a cero cuando $\epsilon \to 0$: un enfriamiento reversible
exige un termostato a cada temperatura intermedia --- el calor nunca ha de atravesar un salto finito.
\end{solution}

\begin{solution}{exo:b1:second-law-entropy:12}
$2^{-10} = 10^{-3}$; $2^{-100} = \num{8e-31}$; $2^{-10^{22}} = 10^{-3\times10^{21}}$.
$k_B\ln2^N = Nk_B\ln2 = nR\ln2$. ``Imposible'' significa en rigor
probabilidad nula; aquí vale $10^{-3\times10^{21}}$ --- no es cero, pero es un número tan pequeño
que jamás podrá observarse diferencia alguna: la irreversibilidad es
estadística a la escala de $10^{23}$.
\end{solution}

\begin{solution}{pb:b1:second-law-entropy:1}
\textbf{1.} $Q = mc(T_0 - T_1) = 0.25 \times 4180 \times (-70) = \qty{-73}{kJ}$
para el café: la habitación gana \qty{73}{kJ}, $\Delta S_{\text{hab}} = 73150/293
= \qty{+250}{J/K}$.

\textbf{2.} $mc\ln(T_0/T_1) = 1045\ln(293/363) = \qty{-224}{J/K}$.

\textbf{3.} $S_{\mathrm{created}} = -224 - (-250) = \qty{+26}{J/K} > 0$.

\textbf{4.} $\Delta S = 1045\ln(330/363) = \qty{-99.6}{J/K}$; $Q = 1045 \times
(-33) = \qty{-34.5}{kJ}$, $S_{\mathrm{exch}} = \qty{-117.7}{J/K}$: creados
\qty{18}{J/K} hasta ahí --- casi toda la irreversibilidad ocurre mientras el
café está más caliente.

\textbf{5.} $\Delta S - Q/T_0 = mc[\ln(T_0/T_1) - (T_0 - T_1)/T_0] = mc[x - 1 -
\ln x]$ con $x = T_1/T_0$; $f'(x) = 1 - 1/x$ se anula en $x = 1$, donde
$f = 0$, y $f'' = 1/x^2 > 0$: mínimo $0$.

\textbf{6.} No: los estados inicial y final son los mismos, el calor es el
mismo y la habitación es el mismo termostato --- la entropía creada la
fijan los estados, no la rapidez.

\textbf{7.} $T_f = \qty{350}{K}$; creados $4180\ln(350^2/120000) = \qty{86}{J/K}$.

\textbf{8.} Reversible: la entropía total se conserva: $mc\ln(T_f'/T_1) + mc\ln(T_f'/
T_2) = 0$, luego $T_f'^2 = T_1T_2$.

\textbf{9.} $T_f' = \sqrt{120000} = \qty{346.4}{K}$; las aguas pierden $mc[(T_1 -
T_f') + (T_2 - T_f')] = 4180 \times (53.6 - 46.4) = \qty{30}{kJ}$, entregados como
trabajo.

\textbf{10.} Caliente: $4180\ln(346.4/400) = \qty{-601}{J/K}$; fría: $4180\ln(346.4/
300) = \qty{+601}{J/K}$: suma nula, como exige la reversibilidad.

\textbf{11.} $T_0S_{\mathrm{created}} = 300 \times 86 = \qty{26}{kJ}$, cerca de los
\qty{30}{kJ} recuperables: la entropía creada, multiplicada por la
temperatura de referencia, es el trabajo que el camino irreversible tiró.

\textbf{12.} Haría falta una máquina que trabajara entre dos aguas tibias
cuyas temperaturas no dejan de cambiar --- unas decenas de kilojulios para
una instalación grande: no compensa, pero el segundo principio dice que el
trabajo estaba ahí.

\textbf{13.} Recibida del interior $Q_c/T_c = 40/278 = \qty{0.144}{W/K}$;
cedida a la habitación $Q_0/T_0 = (40 + W)/293$; balance a lo largo de un ciclo:
$0 = 0.144 - (40 + W)/293 + S_c$.

\textbf{14.} $S_c \geq 0$: $(40 + W)/293 \geq 0.144$, $W \geq 42.2 - 40 =
\qty{2.2}{W}$; COP $\leq 40/2.2 = 18$ ($= T_c/(T_0 - T_c)$).

\textbf{15.} $\mathrm{COP}_{\max} = T_c/(T_0 - T_c)$: para un congelador $255/38 = 6.7$
--- cuanto más fría es la cámara, más caro sale cada julio extraído.

\textbf{16.} Con $W = 15$: $S_c = (55/293) - 0.144 = \qty{0.044}{W/K}$ creados
cada segundo.

\textbf{17.} $W \geq 200 \times 15/278 = \qty{10.8}{W}$; la máquina vierte en la
habitación los \qty{200}{W} más su propio trabajo mientras el interior sigue
frío: la ganancia neta de la habitación es $W$ --- un frigorífico abierto es una estufa.

\textbf{18.} $W = 0$ daría $S_c = Q_c(1/T_0 - 1/T_c) < 0$: prohibido --- el
\omterm{prop:b1:second-law-entropy:clausiuskelvin}{enunciado de Clausius}.

\textbf{19.} Antes: $1$ (todas a la izquierda); después: $2^N$; $\Delta S = k_B\ln2^N =
Nk_B\ln2$.

\textbf{20.} $Nk_B = nR$: $\Delta S = nR\ln2$, el resultado de la \omterm{rem:b1:first-law:irreversible}{expansión de Joule}.

\textbf{21.} Un bit por molécula: $N$ bits; y $\Delta S/(k_B\ln2) = N$ --- la
entropía cuenta la información que haría falta para localizar las
moléculas.

\textbf{22.} $2^{-50} = \num{9e-16}$; para un \omterm{def:b1:kinetic-theory:scales}{mol} $2^{-6\times10^{23}} =
10^{-1.8\times10^{23}}$.

\textbf{23.} Sin calor ($Q = 0$): intercambiada $0$, creada $nR\ln2$ --- el
caso más puro: ni trabajo ni calor, solo la dispersión del gas.

\textbf{24.} $W = nRT\ln2$ recibido, $Q = -nRT\ln2$ cedido a la habitación,
$S_{\mathrm{exch}} = -nR\ln2$ (la entropía del gas vuelve a su valor
inicial, reversiblemente): deshacer la expansión libre cuesta el trabajo
$nRT\ln2$ y vierte $nR\ln2$ de entropía en la habitación --- el universo se
queda con los $nR\ln2$ creados de una vez por todas.

\textbf{25.} La entropía no se conserva en nada real --- solo en el límite
reversible idealizado; el principio prohíbe todo proceso que la destruyera
(calor cuesta arriba, trabajo de un solo termostato, desmezcla gratuita);
y cada brizna creada es trabajo perdido, $T_0S_{\mathrm{created}}$ de él.
\end{solution}
